\documentclass[12pt, ukrainian]{article}
\usepackage[a4paper]{geometry}
\setlength\parindent{0mm}
\geometry{top=1.0cm,bottom=1.0cm,left=1.3cm,right=1.3cm}
\pagestyle{empty}
\usepackage[T2A]{fontenc}
\usepackage[utf8]{inputenc}
\usepackage[ukrainian]{babel}
\usepackage{graphicx}
\usepackage{caption}
\usepackage{caption}
\DeclareCaptionFormat{nocaption}{}
\captionsetup{format=nocaption,aboveskip=0pt,belowskip=0pt}
\usepackage[Export]{adjustbox}
\adjustboxset{max size={0.9\linewidth}{0.9\paperheight}}
\usepackage{verbatim, listings, clrscode3e, fancyvrb}
\def\code#1{\Verb!#1!}
\begin{document}
%begin
\begin {large}

\textbf{\_\_.11.2025 Контрольна робота, варіант  \No { 131 }}\par


{

\begin{center}
\textbf{Завдання 1}
\end{center}
Нехай int var=13; Один потік збільшив змінну var на 3, а інший збільшив її в три рази.

гонка даних (data race), стан гонки (race condition)

Виберіть усі правильні твердження (якщо такі є).\par
\vspace{2mm}
(\,1\,) Тут гарантовано нема гонки даних, але може бути стан гонки.%bad

(\,2\,) Тут гарантовано нема стану гонки, але може бути гонка даних.%bad

(\,3\,) Тут гарантовано є стан гонки, але гонки даних може не бути.%good

(\,4\,) Тут може не бути ані гонки даних, ані стану гонки.%bad



\begin{center}
\textbf{Завдання 2}
\end{center}
Виберіть усі правильні твердження (якщо такі є).\par
\vspace{2mm}

(\,1\,) Критична секція це неперервна ділянка коду програми.%bad

(\,2\,) За різних запусків програми на одних тих самих вхідних даних може виникати хронологічно різна послідовність обчислень.%good

(\,3\,) Об’єкт future можна скопіювати та роздати копії різним потокам, щоб вони змогли побачити результат.%bad

(\,4\,) Щоб розблокувати заблокований м’ютекс \code{recursive\_mutex}, його завжди достатньо розблокувати один раз.%bad



\begin{center}
\textbf{Завдання 3}
\end{center}
Виберіть усі правильні твердження (якщо такі є).\par
\vspace{2mm}

(\,1\,) За використання атомарних змінних структура даних є структурою даних без блокувань.%bad

(\,2\,) Реалізація структури даних без блокувань не може використовувати м’ютекси.%good

(\,3\,) Структура даних без блокувань гарантує завершення виконання довільної операції над цією структурою даних за обмежений час.%bad

(\,4\,) Вільна від блокувань структура даних гарантує, що жоден потік не буде очікувати доступу до виконання операцій над цією структурою даних.%bad



\begin{center}
\textbf{Завдання 4}
\end{center}
Виберіть усі правильні твердження (якщо такі є).\par
\vspace{2mm}

(\,1\,) Автоматичні змінні не можуть розділятися між потоками.%bad

(\,2\,) Потік мови С++ не може отримати доступ до автоматичних змінних іншого потоку.%bad

(\,3\,) Потік може мати доступ до автоматичних змінних іншого потоку, але це небезпечно.%good

(\,4\,) Потік може мати доступ до автоматичних змінних іншого потоку і це безпечно.%bad



\begin{center}
\textbf{Завдання 5}
\end{center}
У цьому фрагменті коду
\begin{verbatim}
template <class T> void f(T a){}

void g(){
    int i=5;
    f(unique_ptr<int>(new int(i)))+(unique_ptr<int>(new int(++i)));
\end{verbatim}
(\,1\,) є невизначена поведінка%good

(\,2\,) є стан гонки%bad

(\,3\,) є гонка даних%bad

(\,4\,) можливий виток ресурсів%good


}
\end {large}


\end{document}

